% =====================================================================
% OLTW: численный пример заполнения двух колонок
% =====================================================================
\begin{frame}{OLTW: численный пример заполнения DP-таблицы}

\textbf{Партитура:} 5 нот гаммы C-мажор $P = (60, 62, 64, 65, 67)$.

\textbf{Поток:} играем первые 4 ноты ровно $O = (60, 62, 64, 65)$.

\textbf{Стоимость:} $c(p_i, o_t) = |p_i - o_t|$ (полутоны).

\bigskip

\textbf{Шаг $t=1$, $o_1 = 60$:}
\[
\begin{array}{c|ccccc}
i & 1 & 2 & 3 & 4 & 5 \\
\hline
p_i & 60 & 62 & 64 & 65 & 67 \\
c(p_i, o_1) & 0 & 2 & 4 & 5 & 7 \\
D[i, 1]     & 0 & 2 & 6 & 11 & 18 \\
\end{array}
\]

(Кумулятивная сумма --- это «вертикальный» шаг $D[i,1] = c(p_i, o_1) + D[i{-}1, 1]$.)

\bigskip

\textbf{Предсказание:} $\hat\imath_1 = \argmin_i D[i, 1] = \rlhi{1}$
(первая нота). 

\end{frame}

\begin{frame}{OLTW: пример заполнения (продолжение)}

\textbf{Шаг $t=2$, $o_2 = 62$.} По рекурсии
$D[i,t] = c(p_i, o_t) + \min(D[i{-}1, t{-}1], D[i{-}1, t], D[i, t{-}1])$:

\[
\begin{array}{c|ccccc}
i & 1 & 2 & 3 & 4 & 5 \\
\hline
c(p_i, o_2) & 2 & 0 & 2 & 3 & 5 \\
D[i, 2]     & 2 & 0 & 2 & 5 & 10 \\
\end{array}
\]

\textbf{Предсказание:} $\hat\imath_2 = \argmin_i D[i, 2] = \rlhi{2}$. 

\bigskip

\textbf{Шаг $t=3$, $o_3 = 64$:}
\[
\begin{array}{c|ccccc}
i & 1 & 2 & 3 & 4 & 5 \\
\hline
c(p_i, o_3) & 4 & 2 & 0 & 1 & 3 \\
D[i, 3]     & 6 & 2 & 0 & 1 & 4 \\
\end{array}
\]
$\hat\imath_3 = \rlhi{3}$.

\bigskip

\textbf{Шаг $t=4$, $o_4 = 65$:} $\hat\imath_4 = \rlhi{4}$
(минимум в строке $i=4$).

\bigskip

\textbf{Итог:} OLTW идеально отслеживает $\hat\imath_t = 1, 2, 3, 4$ ---
ровно по диагонали DP-сетки.\quad
\textbf{Память:} хранилось всего 2 колонки.

\end{frame}
